


% \subsection{Enconding of operations on natural numbers in $\muLJ$ and $\cmuLJ$.}














% Dually, the greatest fixed point operator allow to define infinite data types (called co-data), such as streams, infinite trees, etc. For instance, the type of infinite streams of integers can be represented in linear logic with fixed points by $S= \nu X. \N \times X$, and 



% Another interesting feature of the infinitary proofs is that their computational meaning is more explicit compared to finitary proof systems. For instance, it is clear that the following circular proof computes the concatenation of a list and a stream



% \[
% \small
% \vlderivation{
% \vlin{\rr \mu}{}{\N \seqar N}{
% \vlin{\rr {+^1}}{}{N \seqar \unit + \N}
% {
% \vlin{\id}{}{N \seqar N}{\vlhy{}}
% }
% }
% }
% \vlderivation{
% \vlin{\rr \mu}{}{\Gamma \seqar N}{
% \vlin{\rr {+^1}}{}{\Gamma\seqar \unit + \N}
% {
% \vlhy{\Gamma \seqar \N}
% }
% }
% }







% \[
% \small
% \vlderivation{
% \vlin{\rr \mu}{}{\seqar L}
% {
% \vlin{\rr {+^0}}{}{\seqar\unit + (\N \times L)}
% {
% \vlin{\rr \unit}{}{\seqar\unit}{\vlhy{}}
% }
% }
% }
% \qquad 
% \vlderivation{
% \vlin{\rr \mu}{}{\seqar L}
% {
% \vlin{\rr{+^1}}{}{\seqar \unit +(\N \times L) }
% {
% \vliin{\times}{}{\seqar \N \times L}
% {
% \vltr{\cod{n}}{\seqar \N}{\vlhy{}}{\vlhy{}}{\vlhy{}}
% }
% {
% \vltr{\underline{l}}{\seqar L}{\vlhy{}}{\vlhy{}}{\vlhy{}}
% }
% }
% }
% }
% \]



% \[
% \small
% \vlderivation{
% \vlin{\rr \nu}{}{\seqar S}
% {
% \vliin{\rr \times}{}{\seqar N \times S}
% {
% \vldr{\cod {n_0}}{\seqar N}
% }
% {
% \vlin{\rr \nu}{}{\seqar S}
% {
% \vliin{\rr \times}{}{\seqar N \times S}
% {
% \vldr{\cod {n_1}}{\seqar N}
% }
% {
% \vlin{\rr \nu}{}{\seqar S}
% {
% \vlhy{\vdots}
% }
% }
% }
% }
% }
% }
% \]


% \[
% \small
% \vlderivation{
% \vlin{\lr \nu}{}{L \seqar S \to S}
% {
% \vliin{\lr + }{}{\unit +(N \times (S \to S))  \seqar S \to S}
% {
% \vlin{\lr \unit}{}{\unit \seqar S \to S}
% {
% \vlin{\rr \to }{}{ \seqar S \to  S}
% {
% \vlin{\id }{}{S \seqar S}{\vlhy{}}
% }
% }
% }
% {
% \vlin{\lr \times}{}{N \times (S \to S) \seqar  S \to S}
% {
% \vlin{\rr \to}{}{N, S \to S \seqar  S \to S}
% {
% \vlin{\rr \nu}{}{N, S \to S, S \seqar S}
% {
% \vliin{\rr \times}{}{N, S \to S, S \seqar N \times S}
% {
% \vliin{\lr \to}{}{S \to S, S \seqar S}
% {
% \vlin{\id}{}{S \seqar S}{\vlhy{}}
% }
% {
% \vlin{\id}{}{S \seqar S}{\vlhy{}}
% }
% }
% {
% \vlin{\id}{}{N \seqar N}{\vlhy{}}
% }
% }
% }
% }
% }
% }
% }
% \]



% \[
% \small
% \vlderivation{
% \vlin{\rr \to}{}{L \seqar S \to S}
% {
% \vlin{\lr \mu}{\bullet}{L, S \seqar S}
% {
% \vliin{\lr +}{}{\unit +(N \times L), S \seqar S}
% {
% \vlin{\lr \unit}{}{\unit, S \seqar S}{\vlin{\id}{}{S \seqar S}{\vlhy{}}}
% }
% {
% \vlin{\rr \nu}{}{N \times L, S \seqar S}
% {
% \vlin{\lr \times}{}{N \times L, S \seqar N \times S}
% {
% \vliin{\rr \times}{}{N, L, S \seqar N \times S}
% {
% \vlin{\id}{}{N \seqar N}{\vlhy{}}
% }
% {
% \vlin{\lr \mu  }{\bullet}{L, S \seqar S}
% {
% \vlhy{\vdots}
% }
% }
% }
% }
% }
% }
% }
% }
% \]

% \red{Here cut-reduction for recursion}
% \[
% \small
% \vlderivation{
% \vliin{\cut}{}{\Gamma,\Delta \seqar \tau}
% {
%  \vlin{\Nzero}{}{\seqar \N}{\vlhy{}}
%  }
%  {
%   \vliin{\lr \N}{}{\Gamma, \N \seqar \tau}{\vldr{g(\vec x)}{\Gamma \seqar \sigma}}{\vldr{h(x, \vec x, z)}{\Gamma, \sigma \seqar \sigma}}
%   }
% }
% \qquad \cutred\qquad 
% \vldr{g(\vec x)}{\Gamma \seqar \sigma}
% \]
% \[
% \small
% \vlderivation{
% \vliin{\cut}{}{\Gamma \seqar \tau}
% {
%   \vlin{\Nsucc}{}{\Gamma \seqar \N}{\vlhy{\Gamma \seqar \N}}
%  }
%  {
%   \vliin{\lr \N}{}{\Gamma, \N \seqar \tau}{\vldr{g(\vec x)}{\Gamma \seqar \sigma}}{\vldr{h(x, \vec x, z)}{\Gamma, \sigma \seqar \sigma}}
%   }
% }
% \qquad \cutred\qquad 
% \vlderivation{
% \vliq{\contr}{}{\Gamma \seqar \sigma}
% {
% \vliin{\cut}{}{\Gamma,\Gamma \seqar \sigma}
% {
% \vldr{f(x, \vec x)}{\Gamma \seqar \sigma}
% }
% {
% \vldr{h(x, \vec x, z)}{\Gamma, \sigma \seqar \sigma}
% }
% }
% }
% \]
 % \[
 %    \vlinf{\Nzero}{}{\seqar \N}{}
 %    \quad 
 %    \vlinf{\Nsucc}{}{\Gamma \seqar \N}{\Gamma \seqar \N}
 %    \quad
 %    \vliiinf{\lr \N}{}{\Gamma,\Delta, \N \seqar \tau}{\Gamma \seqar \sigma}{\Gamma, \sigma \seqar \sigma}{\Delta, \sigma \seqar \tau}
 %    % \quad
 %    % \vliinf{\cnd}{}{\Gamma, \N \seqar \tau}{\Gamma \seqar \tau}{\Gamma, \N\seqar \tau}
 %    \]


\section{Weak and strong (co)iteration rules}\label{app:simple-types-with-fixed-points}

% \subsection{Further preliminary content}




% \subsection{Some derivations}


% \begin{enumerate}
%     \item\label{enum:app1}  


%     % \item  \label{enum:app4}  Functoriality in $\cmuLJ$.  The case of the greatest fixed point case is:
%     % \[
%     % \small
%     % \vlderivation{
%     % \vlin{\lr \nu}{\bullet}{\nu X \sigma(X)\seqar \nu X \sigma(X)}{
%     % \vlin{\rr {\nu'}}{}{\sigma(\nu X(\sigma(X))\seqar \nu X \sigma(X)}{
%     % \vliq{\sigma}{}{\sigma(\nu X \sigma(X))\seqar \sigma(\nu X \sigma(X))}{
%     % \vlin{\lr \nu}{\bullet}{\nu X \sigma(X)\seqar \nu X \sigma(X)}{\vlhy{\vdots}}
%     % }
%     % }
%     % }
%     % }
%     % \]

%        \item\label{subsec:encoding-operations-natural-numbers}
%   Encoding of $\lr \N$:
%    \[
% \begin{array}{rcl}
% \lr \N &\dfn&
% \vlderivation{
% \vliin{\lr\mu}{}{\Gamma, \N \seqar \tau}
% {
% \vliin{\lr +}{}{\Gamma, 1 + \sigma \seqar \sigma }
% {
% \vlin{\lr \unit}{}{\Gamma, \unit \seqar \sigma }{\vlhy{\Gamma \seqar \sigma}}
% }
% {
% \vlhy{\Gamma, \sigma \seqar \sigma}
% }
% }
% {
% \vlhy{\Delta, \sigma \seqar  \tau}
% }
% }
% \end{array}
% \]

% \end{enumerate}


Our presentation of the iteration rule (i.e., $\lr\mu$)  and the coiteration rule (i.e., $\rr \nu$) for $\muLJ$ is the most permissive one, in that contexts are allowed to appear in both premises:
 \begin{equation}\label{eqn:strong-2-premises}
      \small
    \vliinf{\lr \mu}{}{\Gamma,\Delta, \mu X \sigma(X)  \seqar \tau}{\Gamma, \sigma(\rho)\seqar \rho}{\Delta, \rho \seqar \tau}
  \qquad
   \vliinf{\rr \nu }{}{ \Delta, \Gamma  \seqar \nu X\sigma(X)}{\Delta \seqar  \tau }{  \Gamma, \tau \seqar \sigma(\tau)}
 \end{equation}

 
 Their cut-reduction rules are as in~\Cref{fig:cut-reduction-mu-nu-muLJ}. The reader should notice that the context $\Gamma$ in both the rules of~\eqref{eqn:strong-2-premises} is contracted during cut-reduction. For this reason, alternative presentations of $\lr\mu $ and $\rr \nu$ without  the context $\Gamma$ have been studied in the literature, especially in the setting of linear logic, where contraction cannot be freely used (see, e.g.,~\cite{BaeldeMiller07})
 
  \begin{equation}\label{eqn:weak-2-premises}
      \small
    \vliinf{\lr \mu}{}{\Delta, \mu X \sigma(X)  \seqar \tau}{ \sigma(\rho)\seqar \rho}{\Delta, \rho \seqar \tau}
  \qquad
   \vliinf{\rr \nu }{}{ \Delta  \seqar \nu X\sigma(X)}{\Delta \seqar  \tau }{   \tau \seqar \sigma(\tau)}
 \end{equation}


 
% \begin{figure*}
%     \centering
%     \[
%     \small
%     \vlderivation{
%     \vliin{\cut}{}{\Sigma, \Gamma,\Delta\seqar \tau}
%     {
%      \vlin{\rr \mu}{}{\Sigma \seqar \mu X \sigma(X)}{
%      \vlhy{\Sigma \seqar \sigma(\mu X \sigma(X))}}
%      }
%      {
%   \vliin{\lr \mu}{}{\Gamma,\Delta, \mu X \sigma(X)  \seqar \tau}
%   {\vlhy{\Gamma, \sigma(\rho)\seqar \rho}}{\vlhy{\Delta, \rho \seqar \tau}}    
%   }
%     }
% \]
% \vspace{0.2cm}
% \[\rightsquigarrow\]
% \[
% \vlderivation{
% \vliq{\contr}{}{\Sigma, \Delta,  \Gamma\seqar \tau}
% {
% \vliin{\cut}{}{\Sigma,  \Delta, \Gamma, \Gamma \seqar \tau}
% {
% \vliin{\cut}{}{\Sigma, \Gamma, \Gamma \seqar \rho}
% {
% \vliin{\cut}{}{\Sigma, \Gamma\seqar \sigma(\rho)}
% {
% \vlhy{\Sigma \seqar \sigma(\mu X \sigma(X))}
% }
% {
% \vliq{\sigma}{}{\Gamma, \sigma(\mu X \sigma(X) ) \seqar \sigma(\rho)}
% {
%  \vliin{\lr \mu}{}{\Gamma, \mu X \sigma(X)  \seqar \rho}
%   {\vlhy{\Gamma, \sigma(\rho)\seqar \rho}}{\vlin{\id}{}{ \rho \seqar \rho}{\vlhy{}}}    
% }
% }
% }
% {\vlhy{\Gamma, \sigma(\rho)\seqar \rho}}
% }
% {
% \vlhy{\Delta, \rho \seqar \rho}
% }
% }
% }
% \]
% \vspace{1.5cm}
% \[
%  \small
% \vlderivation{
% \vliin{\cut}{}{\Delta, \Gamma , \Sigma \seqar \tau }
% {
%  \vliin{\rr \nu }{}{\Delta,  \Gamma \seqar \nu X\sigma(X)}{\vlhy{\Delta \seqar  \rho }}{\vlhy{  \Gamma, \rho \seqar \sigma(\rho)}}
%  }
%  {
%  \vlin{\lr \nu }{}{\Sigma,  \nu X\sigma(X) \seqar \tau}{\vlhy{\Sigma , \sigma(\nu X \sigma(X))\seqar \tau}}
%  }
%  }
% \]
% \vspace{0.2cm}
% \[\rightsquigarrow\]
% \[
% \vlderivation{
% \vliq{\contr}{}{\Delta , \Gamma, \Sigma \seqar \tau}
% {
% \vliin{\cut}{}{\Delta , \Gamma, \Gamma, \Sigma , \Delta \seqar \tau}
% {\vlhy{\Delta \seqar \rho}}
% {
% \vliin{\cut}{}{\Gamma, \Gamma, \Sigma , \rho \seqar \tau}
% {
% \vlhy{  \Gamma, \rho \seqar \sigma(\rho)}
% }
% {
% \vliin{\cut}{}{\Gamma ,  \Sigma, \sigma(\rho) \seqar \tau}
% {
% \vliq{\sigma }{}{\Gamma ,  \sigma(\rho) \seqar \sigma(\nu X\sigma(X))}
% {
% \vliin{\rr \nu }{}{\Gamma , \rho \seqar \nu X\sigma(X)}{\vlin{\id}{}{\rho \seqar  \rho}{\vlhy{} }}{\vlhy{  \Gamma, \rho \seqar \sigma(\rho)}}
% }
% }
% {\vlhy{\Sigma , \sigma(\nu X \sigma(X))\seqar \tau}}
% }
% }
% }
% }
% \]
%     \caption{Cut-reduction rules for $\mu$ and $\nu$ in $\muLJ$.}
%     \label{fig:cut-reduction-mu-nu-muLJ}
% \end{figure*}


Single premise formulations of $\lr \mu$ and $\rr \nu$ for both~\eqref{eqn:strong-2-premises} and~\eqref{eqn:weak-2-premises} have been investigated in the literature as well (see, e.g.,~\cite{Clairambault09}):
\begin{equation}\label{eqn:strong-1-premises}
      \small
    \vlinf{\lr \mu}{}{\Gamma, \mu X \sigma(X)  \seqar \tau}{\Gamma, \sigma(\tau)\seqar \tau}
  \qquad
   \vlinf{\rr \nu }{}{  \Gamma  \seqar \nu X\sigma(X)}{  \Gamma \seqar \sigma(\tau)}
 \end{equation}
 \begin{equation}\label{eqn:weak-1-premises}
      \small
    \vlinf{\lr \mu}{}{ \mu X \sigma(X)  \seqar \tau}{ \sigma(\tau)\seqar \tau}
  \qquad
   \vlinf{\rr \nu }{}{ \tau  \seqar \nu X\sigma(X)}{   \tau \seqar \sigma(\tau)}
 \end{equation}


 The corresponding cut-reduction rules for~\eqref{eqn:weak-2-premises},~\eqref{eqn:strong-1-premises}, and~\eqref{eqn:weak-1-premises} can be easily extracted from the ones in~\Cref{fig:cut-reduction-mu-nu-muLJ}. 

 It is worth mentioning, however, that cut-elimination fails in presence of the rules \eqref{eqn:strong-1-premises} and~\eqref{eqn:weak-1-premises}. By contrast, being essentially endowed with a built-in cut, the rules  \eqref{eqn:strong-2-premises} and~\eqref{eqn:weak-2-premises} allow cut-elimination results. 

Perhaps surprisingly, all  formulations of $\lr \mu$ and $\rr \nu$ are equivalent, meaning they can derive each other  crucially using  the cut rule. 

\begin{proposition}\label{prop:weak-derives-strong}
    In $\LJ$ endowed with the rules $\rr \mu$ and $\lr \nu$ (\Cref{fig:unfolding-rules}), the  rules~\eqref{eqn:strong-2-premises}, \eqref{eqn:weak-2-premises},~\eqref{eqn:strong-1-premises}, and~\eqref{eqn:weak-1-premises} and their cut-reduction are inter-derivable.
\end{proposition}
\begin{proof}
It suffices to show that the rules in~\eqref{eqn:weak-1-premises} can derive the rules in~\eqref{eqn:strong-1-premises}. We only show the case of $\lr \mu$, as the case of $\rr \nu$ is symmetric. W.l.o.g., we will treat  $\Gamma$ in~\eqref{eqn:strong-1-premises} as a single formula. The case where $\Gamma$ is an arbitrary context can be recovered using $\lr \times$ and $\rr \times$. The derivation  is the following:
\[
\vlderivation{
  \vliin{\cut}{}{\Gamma, \mu X. \sigma \seqar \tau}
      {
      \vltr{\der}{\Gamma, \mu X. \sigma \seqar \mu X. (\Gamma \times \sigma)}{\vlhy{\ }}{\vlhy{\ }}{\vlhy{\ }}
      }{
      \vlin{\lr \mu}{}{\mu X. (\Gamma \times \sigma)\seqar \tau}
        {
        \vlin{\lr \times}{}{\Gamma \times \sigma (\tau)\seqar \tau}
          {
             \vlhy{\Gamma, \sigma (\tau)\seqar \tau}
          }
        }
      }
}
\]
where $\der$ is the derivation in~\Cref{fig:deriving-strong-from-weak}. It is tedious but routine to show that the above derivation of $\lr \mu$ allows us to derive the cut-reduction rule for~\eqref{eqn:weak-1-premises}.
\end{proof}

\begin{figure*}
    \centering
   \[
\footnotesize
\vlderivation{
\vlin{\rr \mu}{}{\Gamma, \mu X. \sigma \seqar \mu X. (\Gamma \times \sigma)}
   {
   \vliq{\contr}{}{\Gamma, \mu X. \sigma \seqar \Gamma \times \sigma (\mu X. (\Gamma \times \sigma))}{
   \vliin{\rr \times}{}{\Gamma, \Gamma, \mu X. \sigma \seqar \Gamma \times \sigma (\mu X. (\Gamma \times \sigma))  }
     {
     \vlin{\id}{}{\Gamma \seqar \Gamma}{\vlhy{}}
     }{
     \vliin{\cut}{}{\Gamma, \mu X. \sigma \seqar  \sigma (\mu X. (\Gamma \times \sigma)) }
       {
      \vlin{\lr \mu}{}{\mu X. \sigma \seqar \Gamma \to \sigma (\mu X. (\Gamma \times \sigma))}
        {
        \vlin{\rr \to }{}{\sigma ( \Gamma \to \sigma (\mu X. (\Gamma \times \sigma))) \seqar \Gamma \to \sigma (\mu X. (\Gamma \times \sigma))}
          {
          \vliq{\sigma}{}{\Gamma,\sigma ( \Gamma \to \sigma (\mu X. (\Gamma \times \sigma))) \seqar  \sigma (\mu X. (\Gamma \times \sigma))}
            {
            \vlin{\rr \mu}{}{\Gamma, \Gamma \to \sigma (\mu X. (\Gamma \times \sigma)) \seqar \mu X. (\Gamma \times \sigma)}
              {
              \vliq{\contr}{}{\Gamma, \Gamma \to \sigma (\mu X. (\Gamma \times \sigma)) \seqar \Gamma \times \sigma(\mu X. (\Gamma \times \sigma))}
              {
              \vliin{\rr \times }{}{\Gamma ,\Gamma, \Gamma \to \sigma (\mu X. (\Gamma \times \sigma)) \seqar \Gamma \times \sigma(\mu X. (\Gamma \times \sigma))}
                {
                \vlin{\id}{}{\Gamma \seqar \Gamma}{\vlhy{}}
                }{
                \vliin{\lr \to}{}{\Gamma, \Gamma \to \sigma (\mu X. (\Gamma \times \sigma)) \seqar  \sigma(\mu X. (\Gamma \times \sigma))}
                  {
                   \vlin{\id}{}{\Gamma \seqar \Gamma}{\vlhy{}}
                  }{
                  \vlhy{\sigma (\mu X. (\Gamma \times \sigma)) \seqar  \sigma(\mu X. (\Gamma \times \sigma))}
                  }
                }
              }
              }
            }
          }
        }
       }{
       \vliin{\lr \to}{}{\Gamma \to  \sigma (\mu X. (\Gamma \times \sigma)), \Gamma \seqar \sigma (\mu X. (\Gamma \times \sigma)) }{ \vlin{\id}{}{\Gamma \seqar \Gamma}{\vlhy{}}}{\vlhy{\sigma (\mu X. (\Gamma \times \sigma)) \seqar \sigma (\mu X. (\Gamma \times \sigma))}}
       }
     }
    } 
   }
}
\]
    \caption{Deriving~\eqref{eqn:strong-1-premises} from~\eqref{eqn:weak-1-premises}.}
    \label{fig:deriving-strong-from-weak}
\end{figure*}


% \subsection{(Strong) functoriality in $\muLJ$}\label{subsec:functoriality-in-muLJ}


%     Let $\sigma(X) $ and $\rho(X)$ be (possibly open) types that are positive (and negative, respectively) in $X$. 
%     For a proof $P:\Gamma, \tau \seqar \tau'$ we define $\sigma(P): \Gamma, \sigma(\tau)\seqar \sigma(\tau')$ and $\rho(P): \Gamma, \rho(\tau')\seqar \rho(\tau)$ by simultaneous induction as follows:
%     \begin{itemize}
%     \item If $\sigma(X)=X$ then $\sigma(\der)$ is just $\der$. Notice that it is never the  case that $\rho=X$, as $X$ can only occur negatively in $\rho$.
%     \item If $\sigma$ and  $\rho$ are $\unit$ or some $Y\neq X$ then $\sigma(\der)$ and ${\tau}(\der)$ are defined respectively as follows: 
%     \[
%     \small
%     \vlderivation{
%     \vliq{\wk}{}{\Gamma, \sigma \seqar \sigma}{\vlin{\ax}{}{\sigma \seqar \sigma}{\vlhy{}}}
%     }
%     \qquad 
%     \vlderivation{
%     \vliq{\wk}{}{\Gamma, \rho \seqar \rho}{\vlin{\ax}{}{\rho \seqar \rho}{\vlhy{}}}
%     }
%     \]
%     \item If $\sigma=\sigma_1 \imp \sigma_2$ and $\rho= \rho \imp \rho$ then we define $\sigma(\der)$ and $\rho(\der)$ respectively as follows:
%     \[
%     \small
%     \begin{array}{c}
%     \vlderivation{
%     \vlin{\rr \imp}{}{\Gamma, \sigma(\tau)\seqar \sigma(\tau')  }
%      {
%      \vliq{\contr}{}{\Gamma, \sigma(\tau) , \sigma_1(\tau')\seqar   \sigma_2(\tau')}{
%      \vliin{\lr \imp}{}{\Gamma ,\Gamma, \sigma(\tau) , \sigma_1(\tau') \seqar   \sigma_2(\tau')}
%        {
%       \vldr{\sigma_1(\der)}{\Gamma, \sigma_1(\tau') \seqar \sigma_1(\tau)}
%        }
%        {
%         \vldr{\sigma_2(\der)}{\Gamma, \sigma_2(\tau) \seqar \sigma_2(\tau')}
%        }
%        }
%      }
%     } 
%       \end{array}
%     \]
%  \[
%     \begin{array}{c}
%    \vlderivation{
%     \vlin{\rr \imp}{}{\Gamma, \rho(\tau')\seqar \rho(\tau)  }
%      {
%      \vliq{\contr}{}{\Gamma, \rho(\tau'), \rho_1(\tau)\seqar   \rho_2(\tau)}{
%      \vliin{\lr \imp}{}{\Gamma, \Gamma, \rho(\tau') , \rho_1(\tau) \seqar   \rho_2(\tau)}
%        {
%        \vldr{\rho_1(\der)}{\Gamma, \rho_1(\tau) \seqar \rho_1(\tau')}
%        }
%        {
%         \vldr{\rho_2(\der)}{\Gamma, \rho_2(\tau') \seqar \rho_2(\tau)}
%        }
%        }
%      }
%     }
%     \end{array}
%     \]
%     \item If $\sigma=\sigma_1 \times \sigma_2$ and $\rho= \rho_1 \times \rho_2$  then we define $\sigma(\der)$ and $\rho(\der)$ respectively as follows:
%      \[
%      \small
%     \begin{array}{c}
%    \vlderivation{
%      \vliq{\contr}{}{\Gamma, \sigma(\tau)\seqar \sigma(\tau')  }
%       {
%       \vliin{\rr \times}{}{\Gamma, \Gamma, \sigma(\tau)\seqar\sigma(\tau)  }
%        {
%        \vlin{\lr \times }{}{\Gamma,\sigma(\tau)\seqar \sigma_1(\tau')}
%          {
%          \vldr{\sigma_1(\der)}{\Gamma, \sigma_1(\tau) \seqar \sigma_1(\tau')}
%          }
%        }
%        {
%          \vlin{\lr \times }{}{\Gamma,\sigma(\tau)\seqar \sigma_2(\tau')}
%           {
%              \vldr{\sigma_2(\der)}{\Gamma,  \sigma_2(\tau)\seqar \sigma_2(\tau')}
%           }
%        }
%       }
%     } 
%   \end{array}
%     \]
%       \[
%       \small
%     \begin{array}{c}
%     \vlderivation{
%      \vliq{\contr}{}{\Gamma, \rho(\tau')\seqar \rho(\tau)  }
%       {
%       \vliin{\rr \times}{}{\Gamma, \Gamma, \rho(\tau') \seqar \rho(\tau)   }
%        {
%        \vlin{\lr \times }{}{\Gamma, \rho(\tau') \seqar \rho_1(\tau)}
%          {
%          \vldr{\rho_1(\der)}{\Gamma, \rho_1(\tau') \seqar \rho_1(\tau)}
%          }
%        }
%        {
%          \vlin{\lr \times }{}{\Gamma, \rho(\tau') \seqar \rho_2(\tau)}
%           {
%              \vldr{\rho_2(\der)}{\Gamma,  \rho_2(\tau')\seqar \rho_2(\tau)}
%           }
%        }
%       }
%     } 
%  \end{array}
%     \]
%     \item If $\sigma=\sigma_1 + \sigma_2$ and $\rho= \rho_1 + \rho_2$ then we define $\sigma(\der)$ and $\rho(\der)$ respectively as follows:
%     \[
%     \small
%    \vlderivation{
%     \vliin{\lr +}{}{\Gamma, \sigma(\tau)\seqar  \sigma(\tau') }
%      {
%      \vlin{\rr +}{}{\Gamma, \sigma_1(\tau) \seqar  \sigma(\tau')}
%        {
%        \vldr{\sigma_1(\der)}{\Gamma, \sigma_1(\tau) \seqar \sigma_1(\tau')}
%        }
%      }
%      {
%       \vlin{\rr +}{}{\Gamma, \sigma_2(\tau) \seqar  \sigma(\tau')}
%        {
%         \vldr{\sigma_2(\der)}{\Gamma, \sigma_2(\tau) \seqar \sigma_2(\tau')}
%        }
%      }
%     }
%     \]
%     \[
%     \small
%      \vlderivation{
%     \vliin{\lr +}{}{\Gamma, \rho(\tau') \seqar \rho(\tau)  }
%      {
%      \vlin{\rr +}{}{\Gamma, \rho_1(\tau') \seqar \rho(\tau) }
%        {
%        \vldr{\rho_1(\der)}{\Gamma, \rho_1(\tau') \seqar \rho_1(\tau)}
%        }
%      }
%      {
%       \vlin{\rr +}{}{\Gamma, \rho_2(\tau') \seqar \rho(\tau) }
%        {
%         \vldr{\rho_2(\der)}{\Gamma, \rho_2(\tau') \seqar \rho_2(\tau)}
%        }
%      }
%     }
%     \]
%         % \item if $\sigma$ and/or $\rho$ are formed from smaller types by a connective of $\LJ$, then the corresponding functors are defined as usual. \todo{give in appendix}
%         \item if $\sigma(X) = \mu Y \sigma'(X,Y)$ and $\rho(X) = \mu Y \rho'(X,Y)$  then we define $\sigma(\der)$ and $\rho(\der)$ respectively as follows:
%         \[
%         \small
%         \vlderivation{
%         \vlin{\lr \mu}{}{\Gamma, \sigma(\tau) \seqar \sigma(\tau')}{
%         \vlin{\rr \mu}{}{\Gamma, \sigma'(\tau, \sigma(\tau'))\seqar \sigma(\tau')}{
%         \vldr{\sigma'(P,\sigma(\tau'))}{\Gamma, \sigma'(\tau, \sigma(\tau')) \seqar \sigma'(\tau', \sigma(\tau'))}
%         }
%         }
%         }
%         \]
%          \[
%          \small
%         \vlderivation{
%         \vlin{\lr \mu}{}{\Gamma, \rho(\tau') \seqar \rho(\tau)}{
%         \vlin{\rr \mu}{}{\Gamma, \rho'(\tau', \rho(\tau))\seqar \rho(\tau)}{
%         \vldr{\rho'(P,\rho(\tau))}{\Gamma, \rho'(\tau', \rho(\tau)) \seqar \rho'(\tau, \rho(\tau))}
%         }
%         }
%         }
%         \]
%         where $\sigma'(P,\sigma(\tau'))$ (resp., $\rho'(P,\rho(\tau))$) are obtained from the IH for $\sigma'(P,Y)$ (resp., $\rho'(P,Y)$) under substitution of $\sigma(\tau')$ for $Y$ (resp., $\rho(\tau)$), cf.~Remark~\ref{rem:gen-identity+substitution}.
%         \item   if $\sigma(X) = \nu Y \sigma'(X,Y)$ and $\rho(X) = \nu Y \rho'(X,Y)$  then we define $\sigma(\der)$ and $\rho(\der)$ respectively as follows: 
%             \[
%         \small
%         \vlderivation{
%         \vlin{\rr \nu}{}{\Gamma, \sigma(\tau) \seqar \sigma(\tau')}{
%         \vlin{\lr \nu}{}{\Gamma, \sigma(\tau) \seqar \sigma'(\tau', \sigma(\tau))}{
%         \vldr{\sigma'(P,\sigma(\tau))}{\Gamma, \sigma'(\tau, \sigma(\tau)) \seqar \sigma'(\tau', \sigma(\tau))}
%         }
%         }
%         }
%         \]
%          \[
%          \small
%         \vlderivation{
%         \vlin{\rr \nu}{}{\Gamma, \rho(\tau') \seqar \rho(\tau)}{
%         \vlin{\lr \nu}{}{\Gamma, \rho(\tau')\seqar \rho'(\tau, \rho(\tau'))}{
%         \vldr{\rho'(P,\rho(\tau'))}{\Gamma, \rho'(\tau', \rho(\tau')) \seqar \rho'(\tau, \rho(\tau'))}
%         }
%         }
%         }
%         \]
%         where $\sigma'(P,\sigma(\tau))$ (resp., $\rho'(P,\rho(\tau'))$) are obtained from the IH for $\sigma'(P,Y)$ (resp., $\rho'(P,Y)$) under substitution of $\sigma(\tau)$ for $Y$ (resp., $\rho(\tau')$), cf.~Remark~\ref{rem:gen-identity+substitution}.
%     \end{itemize}
